\documentclass{rvcepaper}
\newcommand{\orrow}{\multicolumn{6}{|c|}{\bfseries OR}\\ \hline}
\newcommand{\instrA}[1]{\begin{seeinstr}\item Answer all questions from Part A. Part A questions should be answered in first three pages of the answer book only.\item Answer FIVE full questions from Part B. In Part B question number 2 is compulsory. Answer any one full question from 3 and 4, 5 and 6, 7 and 8, 9 and 10.#1\end{seeinstr}}
\begin{document}

%%%%%%%%%%%%%%%%%%%%%%%%%%%%%%%%%%%%%%%% MATERIAL SCIENCE FOR ENGINEERS
% ---------------- ME232TB, April/May 2024
\begin{seepaper}
\paperinfo{code={ME232TB},examline={III Semester B.E. Examinations April/May -- 2024},dept={Common to All Branches},title={MATERIAL SCIENCE FOR ENGINEERS},time={03 Hours},maxmarks={100}}
\seehead
\instrA{}
\seepartA
\begin{seeA}
\qaq{1}{1.1}{Ceramic materials are generally having \blank[2cm] type of bonds.}{01}{1}{1}
\qa{1.2}{When the electrical behaviour of the material are dominated by the electronic structure of pure metal, the material can be called as \blank[2cm] semi-conductors.}{01}{1}{1}
\qa{1.3}{Give an example of natural composite.}{01}{1}{1}
\qa{1.4}{$T_i$ is having \blank[2cm] crystal structure at room temperature.}{01}{1}{1}
\qa{1.5}{For the application of extremely high temperatures, \blank[2cm] type of basic materials are used.}{01}{1}{1}
\qa{1.6}{$Cu - Ni$ alloys are an example of \blank[2cm] type of defect.}{01}{1}{1}
\qa{1.7}{Deformations by twins are commonly observed in \blank[2cm] type of metals.}{01}{1}{1}
\qa{1.8}{Calculate the co-ordination number of $BCC$ crystal systems.}{01}{1}{1}
\qa{1.9}{External surfaces are an example of \blank[2cm] type of defects in crystal systems.}{01}{1}{1}
\qa{1.10}{Viscoelastic deformation is observed significantly in \blank[2cm] basic material.}{01}{1}{2}
\qa{1.11}{Diamond pyramid indenter is used in \blank[2cm] hardness testing.}{01}{1}{2}
\qa{1.12}{Addition of chromium (Cr) to steel enhances \blank[2cm] property.}{01}{1}{2}
\qa{1.13}{Which property of dielectric materials is measured in materials testing?}{01}{1}{2}
\qa{1.14}{Photodiodes serves as a \blank[2cm] in optical communication systems.}{01}{1}{2}
\qa{1.15}{Name the hardening process which introduces nitrogen into the surface layer of steel to produce nitrides at the surface layer of the component.}{01}{1}{3}
\qa{1.16}{\blank[2cm] is a critical process in electronics manufacturing primarily used in surface-mount technology.}{01}{1}{3}
\qa{1.17}{Aggregated composites possess \blank[2cm] strength compared to fiber reinforced composites.}{01}{1}{3}
\qa{1.18}{Sol-gel process is an example of \blank[2cm] approach of synthesis of nanomaterials.}{01}{1}{4}
\qa{1.19}{Crystalline structures of materials are measured on the basis of lattice spacing and electron interaction with it, in \blank[2cm] characterization technique.}{01}{1}{4}
\qa{1.20}{Name characterization technique for nanostructure in which atomic structures of the materials are observed using electrons passing through ultra-thin specimens.}{01}{1}{4}
\end{seeA}
\seepart{PART-B}
\begin{seeB}
\qb{2}{a}{Describe the following bonds for solid materials with examples.
\begin{rom}\item Ionic bonds\item Secondary bonds and\item Metallic bonds.\end{rom}}{08}{2}{1}
 & b & Explain the following crystallographic features of a solid material.
\begin{rom}\item $FCC$ crystal structure with example\item Energy bands in insulator\item Edge dislocation in a metal.\end{rom} & 08 & 2 & 1\\ \hline
\qb{3}{a}{Define and explain the following thermal properties of a material.
\begin{rom}\item Thermo-electric effect\item Linear thermal expansion coefficient\item Thermal shock.\end{rom}}{08}{2}{2}
 & b & Explain the principle of working of following materials.
\begin{rom}\item Piezoelectric materials and\item Thermocouples.\end{rom} & 08 & 3 & 2\\ \hline
\orrow
\qb{4}{a}{With a stress-strain curve, explain the following properties of a material.
\begin{rom}\item Young's Modulus\item Ultimate Tensile Strength and\item Toughness.\end{rom}}{08}{3}{2}
 & b & Define the following mechanical properties of a solid material.
\begin{rom}\item Fracture toughness\item Hardness\item Fatigue limit.\end{rom} & 08 & 2 & 2\\ \hline
\qb{5}{a}{Explain the properties and applications of stainless steel.}{06}{2}{2}
 & b & Mention the dopant used for p-type and n-type semiconductors. & 04 & 2 & 2\\
 & c & Explain the properties and applications of ceramics. & 06 & 2 & 2\\ \hline
\orrow
\qb{6}{a}{Explain the properties and applications of fiber reinforced composites.}{08}{2}{2}
 & b & Classify forming processes of structural materials and explain any two forming processes. & 08 & 2 & 2\\ \hline
\qb{7}{a}{Explain the steps involved in thermal oxidation method of semiconductor devices.}{08}{3}{3}
 & b & Describe the cause, effect and prevention methods of any two defects of heat treatment process of metals. & 08 & 2 & 3\\ \hline
\orrow
\qb{8}{a}{Differentiate between annealing and normalizing process of ferrous materials.}{08}{3}{3}
 & b & Describe hardening of steel with the help of $TTT$ diagram. & 08 & 2 & 3\\ \hline
\qb{9}{a}{Define nanomaterials and describe the applications of nanomaterials with example.}{08}{2}{4}
 & b & Describe the effect of particle size and surface area of nanomaterials on the mechanical and optical properties of it. & 08 & 2 & 4\\ \hline
\orrow
\qb{10}{a}{Describe the significant properties and applications of carbon nanotubes.}{08}{2}{4}
 & b & Describe the principle of working of $X$-Ray Diffraction ($XRD$) techniques for the characterization of solid materials. & 08 & 3 & 4\\ \hline
\end{seeB}
\end{seepaper}

\end{document}
